\subsection{Data Generation Pipeline}

Our pipeline has four stages. In the first stage, we generate
time-domain waveforms using the IMRPhenomD phenomenological model
\citep{Khan2016} as implemented in PyCBC. Each waveform is sampled at
$\Delta t = 1/4096$\,s with a lower frequency cutoff of
$f_{\rm low} = \SI{20}{Hz}$. For the injected studies we fix
$m_1 = 30 M_\odot$ and $m_2 = 25 M_\odot$ unless otherwise stated. In
the second stage, we apply one of the three modulations defined in
Eqs.~\eqref{eq:amp_dev}--\eqref{eq:freq_dev} to the GR waveform, with
$\beta$ drawn uniformly from one of several predefined ranges. In the
third stage, we generate noise either from the analytic aLIGO
Zero-Detuned-High-Power PSD \citep{psd_model} or from real H1 strain
segments retrieved from the Gravitational-Wave Open Science Center
(GWOSC). The signal is scaled by a single fixed factor, computed once
from a GR reference waveform, and added to a random noise segment. In
the fourth stage, every sample is whitened with the empirical PSD,
bandpass filtered between $20$ and $\SI{500}{Hz}$, windowed to a
$\SI{4}{s}$ region centered on the merger peak, amplitude-normalized
to unit variance (for the amplitude-controlled studies), and clipped
at $\pm 100\sigma$.

\subsection{GW150914 Template Study}

For the GW150914 study we retrieve the raw strain file
\texttt{H-H1\_GWOSC\_4KHZ\_R1-1126259447-32.hdf5} from the GWOSC
event API for GWTC-1-confident/GW150914/v3. The file contains
$131{,}072$ samples at $\SI{4096}{Hz}$ (32 seconds), starting at
GPS $1126259447$. We estimate the noise PSD using the first 4 seconds
of the file with Welch's method (segment length $4096$ samples). The
full 32-second segment is whitened with this PSD and bandpass filtered
between $20$ and $\SI{500}{Hz}$. The GW150914 event at GPS
$1126259462.4$ is located at sample $63078$ of the file. We extract a
$\SI{4}{s}$ window centered on this position, giving the GR template
$h_{\rm GR}(t)$. The modified template $h_{\rm mod}(t)$ is obtained
by applying the amplitude modulation of Eq.~\eqref{eq:amp_dev}. For
each training sample, we draw a random 4-second noise segment from
the same file, excluding the event region, add a scaled version of
either $h_{\rm GR}$ or $h_{\rm mod}$, normalize to unit variance, and
clip. This produces controlled counterfactual waveforms from the
GW150914-derived template; it does not modify the historical detector
observation.

\subsection{Dataset Splitting}

We use a stratified split of $70\%$ training, $15\%$ validation and
$15\%$ test\_seen, together with a separate held-out test\_unseen set.
The test\_unseen set contains only the deviated class for one specific
$\beta$ range (type C), which is excluded from training. This measures
generalization to unseen deviation strengths.

\subsection{Training Protocol}

We use Adam with learning rate $3 \times 10^{-3}$, batch size $32$ or
$64$, cosine annealing of the learning rate, gradient clipping at
norm $1$, and early stopping on validation loss with patience
$20$--$50$ epochs. Training and evaluation are performed on CPU using
PyTorch. Each experiment takes one to five minutes.

\subsection{Baseline Classifier}

For comparison we train a class-balanced random forest with $300$ trees
on nine of the ten hand-crafted features. This provides a non-neural
baseline and allows us to identify the most informative features.
